% De Officiis --- headnote
% de-officiis, 44 BC, Puteoli / Tusculum (dedicated to his son Marcus)

\textit{\textit{De Officiis} --- On Duties --- is Cicero's last completed
prose work, written in the autumn of 44 BC in the months after the
assassination of Caesar, while Mark Antony's power was rising and the
proscriptions that would cost Cicero his life were less than a year
away. He cast it as a long letter of instruction to his only son
Marcus, then twenty and studying philosophy at Athens under the
Peripatetic Cratippus --- a son whose record was uneven and whose
father's anxiety shows through the work's affectionate, admonitory
tone. For the first two books Cicero follows the treatise of the Stoic
Panaetius on appropriate action (\textit{peri tou kathēkontos}); the
third part, which Panaetius announced but never wrote, Cicero supplies
himself, and says so plainly. The subject is \textit{officium} ---
duty, the fitting act --- not as abstract theory but as the practical
question of how a decent man should conduct himself in a commonwealth
coming apart.}

\textit{The work falls into three books on a single rising argument.
The first treats the honorable (\textit{honestum}) and traces it to
four sources: the perception of truth, which yields wisdom; the
preservation of human fellowship, which yields justice and generosity;
the greatness and strength of an unconquered spirit, which yields
courage; and order and measure in word and deed, which yield the
seemly --- the \textit{decorum} that is Cicero's distinctive
contribution to ethics. The second turns to the expedient
(\textit{utile}): how goodwill, trust, and glory are truly won, with
the insistence --- against the cynic --- that nothing genuinely
advantageous is ever at odds with the honorable. The third confronts
the cases where advantage seems to clash with honor and dissolves the
conflict through a gallery of tests: the grain-merchant sailing to
famine-struck Rhodes, the ring of Gyges that confers invisibility, and
above all Regulus, the consul who kept his oath to Carthage though it
meant returning to torture and death.}

\textit{No other work of Cicero's has had so long an afterlife. Its
Roman, example-driven ethics --- duty argued through Fabricius and
Regulus rather than through syllogism --- made it a schoolbook for
seventeen centuries, a favorite of the Fathers, of the humanists, and
of Frederick the Great, and among the first books to be set in type
after the Bible. Yet beneath the public teaching runs a private
urgency: a father writing in haste and in danger, handing his son the
distilled conviction of a lifetime in the Forum --- that the honorable
and the expedient are finally one, and that a man who betrays the first
to grasp the second has miscalculated even his own advantage. Within
the year Cicero would be dead at the hands of the men this book quietly
judges; it stands as something close to his testament.}
